% =============================================================================
%  Cup Manipulator: Argument Modes, IK/Kinematics, and Pulley Cable Routing
%  Covers: script_cup_manipulator_pydrake.py  |  test_drive_pulley.py
% =============================================================================

\documentclass[11pt, a4paper]{article}

\usepackage{amsmath, amssymb, bm}
\usepackage{geometry}
\usepackage{listings}
\usepackage{xcolor}
\usepackage{booktabs}
\usepackage{hyperref}
\usepackage[backgroundcolor=blue!5, linecolor=blue!40, linewidth=1pt,
            innertopmargin=6pt, innerbottommargin=6pt,
            innerleftmargin=7pt, innerrightmargin=7pt]{mdframed}

\geometry{margin=2.2cm}

\lstset{
  language=Python,
  basicstyle=\ttfamily\small,
  keywordstyle=\color{blue!70!black},
  commentstyle=\color{gray},
  stringstyle=\color{orange!80!black},
  frame=single,
  breaklines=true,
  columns=fullflexible,
  backgroundcolor=\color{gray!8},
  xleftmargin=6pt, xrightmargin=6pt,
  aboveskip=4pt, belowskip=4pt,
}

\newcommand{\bq}{\bm{q}}
\newcommand{\bp}{\bm{p}}

% =============================================================================
\title{\textbf{Cup Manipulator: Modes, Kinematics \& Pulley Routing}\\
       \large \texttt{script\_cup\_manipulator\_pydrake.py} \;|\; \texttt{test\_drive\_pulley.py}}
\date{March 2026}
\author{Isaac-sim robotics notes}
% =============================================================================

\begin{document}
\maketitle

% ---------------------------------------------------------------------------
\section{Script Overview: \texttt{script\_cup\_manipulator\_pydrake.py}}
% ---------------------------------------------------------------------------

The script drives a 2-DOF revolute cup manipulator (joints
\texttt{link1\_base} $q_1$, \texttt{link2\_link1} $q_2$) with an optional
3-D pendulum ball attached to link~2.  It is invoked via \verb|--mode|:

\begin{center}
\begin{tabular}{ll}
\toprule
\textbf{Mode} & \textbf{Description} \\
\midrule
\texttt{scene-viz}    & Static Meshcat display, no physics \\
\texttt{simulation}   & Full PhysX dynamics, PD torque control \\
\texttt{joint-motion} & Sinusoidal joint sweep, excites ball swing \\
\texttt{run-all-jts}  & Interactive terminal: type angles, robot moves \\
\texttt{min-jerk-joint} & Smooth point-to-point with minimum-jerk profile \\
\bottomrule
\end{tabular}
\end{center}

% ---------------------------------------------------------------------------
\section{Forward Kinematics (2-DOF Planar)}
% ---------------------------------------------------------------------------

With link lengths $L_1 = 342.47\,\text{mm}$ and $L_2 = 190.00\,\text{mm}$
(both lying in the horizontal $XY$ plane at $Z = \text{const}$):

\begin{equation}
  \bp_{ee} =
  \begin{pmatrix} x_{ee} \\ y_{ee} \end{pmatrix}
  =
  \begin{pmatrix}
    L_1 \cos q_1 + L_2 \cos(q_1 + q_2) \\
    L_1 \sin q_1 + L_2 \sin(q_1 + q_2)
  \end{pmatrix}
  \label{eq:fk}
\end{equation}

\noindent Drake computes this implicitly; the code reads the EE position via:

\begin{lstlisting}
cup_body = plant.GetBodyByName("link2", model_instance)
X_WC = plant.CalcRelativeTransform(context,
                                    plant.world_frame(),
                                    cup_body.body_frame())
p_ee = X_WC.translation()   # [x, y, z] world frame
\end{lstlisting}

% ---------------------------------------------------------------------------
\section{Minimum-Jerk Joint-Space Profile (\texttt{min-jerk-joint} mode)}
% ---------------------------------------------------------------------------

Given a move of duration $T$ and normalised time $s = t/T \in [0,1]$,
the minimum-jerk velocity profile is:

\begin{equation}
  h(s) = 10s^3 - 15s^4 + 6s^5
  \label{eq:minjerk-pos}
\end{equation}

\begin{equation}
  \dot{h}(s) = \frac{30s^2 - 60s^3 + 30s^4}{T}, \qquad
  \ddot{h}(s) = \frac{60s - 180s^2 + 120s^3}{T^2}
  \label{eq:minjerk-vel}
\end{equation}

$h(0)=0$, $h(1)=1$; first and second derivatives vanish at both endpoints,
giving smooth ($C^2$) motion with zero jerk at start and end.

\begin{lstlisting}
def min_jerk_profile(t: float, T: float):
    s = np.clip(t / T, 0.0, 1.0)
    h     = 10*s**3 - 15*s**4 + 6*s**5
    hdot  = (30*s**2 - 60*s**3 + 30*s**4) / T
    hddot = (60*s  - 180*s**2 + 120*s**3) / T**2
    return h, hdot, hddot

# Apply to joint space
q_cmd = q_start + h * (q_target - q_start)
\end{lstlisting}

% ---------------------------------------------------------------------------
\section{Script Overview: \texttt{test\_drive\_pulley.py}}
% ---------------------------------------------------------------------------

Visualises the tendon/cable routing of the cable-driven manipulator in
Meshcat.  The cable transmits torque from the drive pulley on body $q_1$
(\texttt{pulley\_htd\_5m\_60t}) through two idler bearings (623ZZ) to the
large driven pulley on body $q_2$ (\texttt{link2\_tendon}).

\medskip
\begin{mdframed}
\textbf{Cable route (ordered)}:\\
Start anchor $\xrightarrow{\;}$
Drive pulley (HTD 5M 60T, $r_p \approx 47.75\,\text{mm}$) $\xrightarrow{\;}$
Idler-L \;/\; Idler-R (623ZZ) $\xrightarrow{\;}$
Big driven pulley ($q_2$ body) $\xrightarrow{\;}$
End anchor
\end{mdframed}

% ---------------------------------------------------------------------------
\section{URDF Visual-Origin Transform}
% ---------------------------------------------------------------------------

Every pulley centroid is computed from its OBJ mesh centre
$\bm{p}_{\text{local}}$ and the URDF \texttt{<visual><origin>} tag:

\begin{equation}
  \bm{p}_{\text{body}}
  = R_z(\psi)\,R_y(\theta)\,R_x(\phi)\;\bm{p}_{\text{local}}
  \;+\; \bm{r}_{\text{xyz}}
  \label{eq:vis-origin}
\end{equation}

where $(\phi,\theta,\psi)$ are roll-pitch-yaw from the URDF \texttt{rpy}
attribute.

\begin{lstlisting}
# PulleyBase._compute_centroid()
roll, pitch, yaw = self.vis_rpy
Rx = np.array([[1,0,0],[0,np.cos(roll),-np.sin(roll)],
                        [0,np.sin(roll), np.cos(roll)]])
Ry = ...   # standard Ry(pitch)
Rz = ...   # standard Rz(yaw)
R  = Rz @ Ry @ Rx

centroid_body = R @ centroid_local + np.array(self.vis_xyz)
\end{lstlisting}

% ---------------------------------------------------------------------------
\section{External Tangent Between Two Pulleys}
% ---------------------------------------------------------------------------

To route the cable between two circular pulleys with centres $\bm{c}_1,
\bm{c}_2$ (world frame) and radii $r_1, r_2$, the external tangent contact
points $T_1, T_2$ satisfy:

\begin{equation}
  D = \|\bm{c}_2 - \bm{c}_1\|, \qquad
  \hat{\bm{d}} = \frac{\bm{c}_2 - \bm{c}_1}{D},\qquad
  \hat{\bm{n}}_\perp = (-\hat{d}_y,\; \hat{d}_x)
\end{equation}

\begin{equation}
  \cos\alpha = \frac{r_1 - r_2}{D},\qquad
  \sin\alpha = \sqrt{1 - \cos^2\!\alpha}
\end{equation}

\begin{equation}
  \bm{n} = \cos\alpha\;\hat{\bm{d}} + \epsilon\,\sin\alpha\;\hat{\bm{n}}_\perp,
  \quad \epsilon \in \{+1,-1\}
\end{equation}

\begin{equation}
  T_1 = \bm{c}_1 + r_1\,\bm{n},\qquad
  T_2 = \bm{c}_2 + r_2\,\bm{n}
  \label{eq:tangent}
\end{equation}

$\epsilon$ selects the upper/lower branch.
For an \emph{internal} tangent (cable crosses between pulleys) replace
$\cos\alpha = (r_1+r_2)/D$ and negate $r_2$ in $T_2$.

\begin{lstlisting}
# PulleyBase.compute_tangent()  (works in XY; Z preserved)
d     = p2 - p1                          # 2-D displacement
D     = np.linalg.norm(d)
d_hat = d / D
perp  = np.array([-d_hat[1], d_hat[0]])  # CCW normal

cos_a = (r1 - r2) / D                   # external tangent
sin_a = np.sqrt(max(0.0, 1.0 - cos_a**2))
n = cos_a * d_hat + branch * sin_a * perp   # branch = +/-1

T1 = np.array([p1[0] + r1*n[0], p1[1] + r1*n[1], c1[2]])
T2 = np.array([p2[0] + r2*n[0], p2[1] + r2*n[1], c2[2]])
\end{lstlisting}

After computing tangents in the world frame, Drake FK maps them
back into the respective body frames for cable drawing:

\begin{lstlisting}
# _Xw returns (R, t) for a body in the world frame
R1, t1 = _Xw(cfg1.body_name)
c1_w   = R1 @ cfg1.centroid + t1        # centroid in world

T1_w, T2_w = PulleyBase.compute_tangent(
    c1_w, cfg1.radius, c2_w, cfg2.radius, branch)

# Back to body frame (for FK consistency in draw_cables)
T1_body = R1.T @ (T1_w - t1)
T2_body = R2.T @ (T2_w - t2)
\end{lstlisting}

% ---------------------------------------------------------------------------
\section{Pulley Pitch Radius (HTD 5M 60T)}
% ---------------------------------------------------------------------------

Both the drive and driven pulleys share the same HTD 5M 60T belt spec:

\begin{equation}
  r_p = \frac{N \cdot p}{2\pi}
      = \frac{60 \times 5\,\text{mm}}{2\pi}
      \approx 47.75\,\text{mm}
  \label{eq:pitch-radius}
\end{equation}

\begin{lstlisting}
belt_teeth   = 60
belt_pitch_m = 0.005          # 5 mm pitch (HTD-5M)
pitch_radius = belt_teeth * belt_pitch_m / (2 * np.pi)  # ~0.04775 m
\end{lstlisting}

\end{document}
